\cotefile{77}

\scanpage{98}

Considérons alors

\[
t'=\tilde f(x\unsure{+|-}iy)=x\wedge y/q(x)=x\wedge y
\]

il est orth. à $F$, \inserted{et partout non nul} donc $t$ de la forme $t=\lambda t'$, de plus $q(t)=q(\lambda t')=-1$, donc $\lambda^2=1$. \deleted{Considérons}

\deleted{Donc \unsure{alors}}

\[
\left\{\begin{aligned} &q(x)=q(\lambda y)=1,\quad\varphi(x,\lambda y)=0\\ &\deleted{\unsure{x+iy|x+i\lambda y}}\tilde f(x+i\lambda y)=x\wedge\lambda y=\lambda x\wedge y=\lambda t'=t \qquad\text{O.K.}\\ &\deleted{f(x+i\lambda y)=}\end{aligned}\right.
\]

\uline{Corollaire} En inversant l'iso. $f$, on trouve une inclusion ouverte (en car. rés. $\ne 2$)

\[
\Sigma\lhook\joinrel\xrightarrow{g}X^{S[i]}\tag*{(24)}
\]

\deleted{\uline{Remarque} Soit $q'=-q$, on a}

\[
\deleted{\delta(q')=-\delta(q)=\omega^{\otimes(-2)}}
\]

\deleted{donc $(E,q',\omega)$ est \deleted{\unsure{anti}}\inserted{structure} spéciale-orthogonale au sens modifié, forme bil. associée}

\[
\deleted{\varphi'(x,y)=-\varphi(x,y)}
\]

\deleted{Si on calque la définition de \inserted{$[\,,\,]=$}$[\,,\,]_q$ pour $q'$, on trouve $[\,,\,]_{q'}=[\unsure{{}^{\prime}},\,]'$ défini par}

\[
\deleted{\varphi'([x,y]',z)\,\omega=2x\wedge y\wedge z}
\]

\deleted{i.e.}

\[
\deleted{\varphi(\unsure{{}^{\prime\prime}}[x,y],z)\,\omega=-\varphi'([x,y],z)\,\omega}
\]

\[
\deleted{[x,y]'=-[x,y].}
\]

\deleted{La formule (21) devient alors}

\[
\deleted{q'([x,y]')=4q'(x)q'(y)\underbrace{\sin^2(\widehat{x,y}\,')}_{1-\deleted{\unsure{4}}\varphi'(x,y)^2/4q'(x)q'(y)}}\tag*{(25)}
\]

% archive sheet number: 97
